\documentclass[12pt,a4paper]{article}

\usepackage[utf8]{inputenc}
\usepackage[T2A]{fontenc}
\usepackage[russian]{babel}
\usepackage{amsmath, amsfonts, amssymb, amsthm}
\usepackage{geometry}
\usepackage{tikz}
\usepackage{hyperref}
\geometry{left=2cm, right=2cm, top=2cm, bottom=2cm}

\newtheorem{theorem}{Теорема}
\newtheorem{lemma}{Лемма}
\theoremstyle{definition}
\newtheorem{definition}{Определение}
\newtheorem{example}{Пример}
\newtheorem{property}{Свойство}

\title{Билет №90 \\ \small Преобразование Фурье. Связь между дискретным и аналоговым преобразованиями Фурье. Свойства преобразования Фурье. Алгоритм быстрого преобразования Фурье. (СпБПУ)}
\author{Сериков Владислав Александрович}
\date{16 мая 2026}

\begin{document}

\maketitle
\tableofcontents
\newpage

\section{Непрерывное преобразование Фурье (НПФ)}

Преобразование Фурье является важнейшим инструментом гармонического анализа, позволяющим переходить от временного представления сигнала к его частотному (спектральному) представлению.

\begin{definition}
Пусть $x(t)$ — абсолютно интегрируемая функция на всей числовой оси, то есть $x(t) \in L_1(\mathbb{R})$. \textbf{Прямым непрерывным преобразованием Фурье} функции $x(t)$ называется функция комплексного переменного $X(\omega)$, задаваемая интегралом:
\begin{equation}
    X(\omega) = \int_{-\infty}^{\infty} x(t) e^{-j\omega t} \, dt,
\end{equation}
где $j$ — мнимая единица ($j^2 = -1$), $\omega$ — круговая частота.
\end{definition}

\begin{definition}
\textbf{Обратным непрерывным преобразованием Фурье} (ОПФ) называется интеграл:
\begin{equation}
    x(t) = \frac{1}{2\pi} \int_{-\infty}^{\infty} X(\omega) e^{j\omega t} \, d\omega.
\end{equation}
\end{definition}

\section{Свойства преобразования Фурье}

Рассмотрим основные свойства преобразования Фурье. Обозначим связь между функцией и её спектром как $x(t) \leftrightarrow X(\omega)$.

\begin{property}[\textbf{Линейность}]
Для любых констант $a, b \in \mathbb{C}$:
$$ a x_1(t) + b x_2(t) \leftrightarrow a X_1(\omega) + b X_2(\omega). $$
\end{property}

\begin{property}[\textbf{Запаздывание / Сдвиг во времени}]
$$ x(t - t_0) \leftrightarrow X(\omega) e^{-j\omega t_0}. $$
\end{property}

\begin{property}[\textbf{Изменение масштаба времени}]
$$ x(at) \leftrightarrow \frac{1}{|a|} X\left(\frac{\omega}{a}\right), \quad a \neq 0. $$
\end{property}

\begin{theorem}[\textbf{О свертке во временной области}]
Пусть функции $x(t)$ и $y(t)$ имеют спектры $X(\omega)$ и $Y(\omega)$ соответственно. Тогда спектр свертки этих функций равен произведению их спектров:
$$ x(t) * y(t) = \int_{-\infty}^{\infty} x(\tau) y(t - \tau) \, d\tau \leftrightarrow X(\omega)Y(\omega). $$
\end{theorem}

\begin{proof}
Запишем определение преобразования Фурье для свертки:
$$ \mathcal{F}\{x(t) * y(t)\} = \int_{-\infty}^{\infty} \left( \int_{-\infty}^{\infty} x(\tau) y(t - \tau) \, d\tau \right) e^{-j\omega t} \, dt. $$
Предполагая, что функции достаточно хорошие (например, из $L_1$), воспользуемся теоремой Фубини и поменяем порядок интегрирования:
$$ \mathcal{F}\{x(t) * y(t)\} = \int_{-\infty}^{\infty} x(\tau) \left( \int_{-\infty}^{\infty} y(t - \tau) e^{-j\omega t} \, dt \right) d\tau. $$
Сделаем замену переменной во внутреннем интеграле: $u = t - \tau$, тогда $dt = du$, а $t = u + \tau$:
$$ \int_{-\infty}^{\infty} y(t - \tau) e^{-j\omega t} \, dt = \int_{-\infty}^{\infty} y(u) e^{-j\omega (u + \tau)} \, du = e^{-j\omega \tau} \int_{-\infty}^{\infty} y(u) e^{-j\omega u} \, du = e^{-j\omega \tau} Y(\omega). $$
Подставим это обратно во внешний интеграл:
$$ \mathcal{F}\{x(t) * y(t)\} = \int_{-\infty}^{\infty} x(\tau) \left( e^{-j\omega \tau} Y(\omega) \right) d\tau = Y(\omega) \int_{-\infty}^{\infty} x(\tau) e^{-j\omega \tau} \, d\tau = Y(\omega) X(\omega). $$
Теорема доказана.
\end{proof}

\begin{theorem}[\textbf{Теорема Парсеваля / Равенство Планшереля}]
Для функций $x(t) \in L_1(\mathbb{R}) \cap L_2(\mathbb{R})$ выполняется равенство энергий во временной и частотной областях:
$$ \int_{-\infty}^{\infty} |x(t)|^2 \, dt = \frac{1}{2\pi} \int_{-\infty}^{\infty} |X(\omega)|^2 \, d\omega. $$
\end{theorem}

\begin{proof}
Представим $|x(t)|^2$ как $x(t)\overline{x(t)}$, где черта означает комплексное сопряжение.
$$ E = \int_{-\infty}^{\infty} x(t) \overline{x(t)} \, dt. $$
Выразим $x(t)$ через обратное преобразование Фурье:
$$ x(t) = \frac{1}{2\pi} \int_{-\infty}^{\infty} X(\omega) e^{j\omega t} \, d\omega. $$
Подставим это выражение в интеграл:
$$ E = \int_{-\infty}^{\infty} \left( \frac{1}{2\pi} \int_{-\infty}^{\infty} X(\omega) e^{j\omega t} \, d\omega \right) \overline{x(t)} \, dt. $$
Меняем порядок интегрирования (по теореме Фубини):
$$ E = \frac{1}{2\pi} \int_{-\infty}^{\infty} X(\omega) \left( \int_{-\infty}^{\infty} \overline{x(t)} e^{j\omega t} \, dt \right) d\omega. $$
Заметим, что внутренний интеграл является комплексным сопряжением прямого преобразования Фурье:
$$ \int_{-\infty}^{\infty} \overline{x(t)} e^{j\omega t} \, dt = \overline{\int_{-\infty}^{\infty} x(t) e^{-j\omega t} \, dt} = \overline{X(\omega)}. $$
Тогда:
$$ E = \frac{1}{2\pi} \int_{-\infty}^{\infty} X(\omega) \overline{X(\omega)} \, d\omega = \frac{1}{2\pi} \int_{-\infty}^{\infty} |X(\omega)|^2 \, d\omega. $$
Теорема доказана.
\end{proof}

\section{Связь между аналоговым и дискретным преобразованиями}

При переходе к цифровой обработке сигналов необходимо перейти от непрерывного времени к дискретному. Этот процесс называется \textbf{дискретизацией} (сэмплированием).

\subsection{Дискретизация сигнала}
Пусть $x(t)$ — непрерывный сигнал. Дискретный сигнал $x[n]$ получается путем взятия отсчетов с периодом $T_s$:
$ x[n] = x(nT_s). $
Математически идеальную дискретизацию удобно описывать умножением непрерывного сигнала на \textit{гребенку Дирака}:
$$ x_s(t) = x(t) \sum_{n=-\infty}^{\infty} \delta(t - nT_s) = \sum_{n=-\infty}^{\infty} x(nT_s) \delta(t - nT_s). $$

\subsection{Спектр дискретизированного сигнала}
Применяя непрерывное преобразование Фурье к $x_s(t)$, получаем:
$$ X_s(\omega) = \mathcal{F}\{x_s(t)\} = \int_{-\infty}^{\infty} \sum_{n=-\infty}^{\infty} x(nT_s) \delta(t - nT_s) e^{-j\omega t} \, dt = \sum_{n=-\infty}^{\infty} x[n] e^{-j\omega nT_s}. $$
Эта формула задает \textbf{Дискретное во времени преобразование Фурье (ДВПФ / DTFT)}.
Используя свойства гребенки Дирака, спектр дискретизированного сигнала также можно выразить через спектр исходного непрерывного сигнала $X(\omega)$:
$$ X_s(\omega) = \frac{1}{T_s} \sum_{k=-\infty}^{\infty} X\left(\omega - \frac{2\pi k}{T_s}\right). $$
\textbf{Связь:} Спектр дискретизированного сигнала является периодическим продолжением спектра аналогового сигнала с периодом $\omega_s = \frac{2\pi}{T_s}$. Если спектр аналогового сигнала не ограничен частотой $\omega_s/2$, происходит наложение спектров (алиасинг). Это суть теоремы Котельникова (Найквиста-Шеннона).

\subsection{Дискретное преобразование Фурье (ДПФ / DFT)}
Поскольку ДВПФ $X_s(\omega)$ является непрерывной функцией от $\omega$, для вычислений на ЭВМ необходимо дискретизировать и спектр. Возьмем $N$ отсчетов спектра на периоде $\omega_s$:
$\omega_k = \frac{\omega_s}{N} k = \frac{2\pi k}{N T_s}$, где $k = 0, 1, \dots, N-1$.
Тогда:
$$ X[k] = \sum_{n=0}^{N-1} x[n] e^{-j \frac{2\pi}{N} kn}. $$
Это и есть формула \textbf{Дискретного преобразования Фурье (ДПФ)}. Обозначим $W_N = e^{-j\frac{2\pi}{N}}$ — поворачивающий множитель.
Тогда ДПФ и ОДПФ записываются как:
$$ X[k] = \sum_{n=0}^{N-1} x[n] W_N^{kn}, \quad x[n] = \frac{1}{N} \sum_{k=0}^{N-1} X[k] W_N^{-kn}. $$

\textbf{Итог связи:} Аналоговое преобразование Фурье $\rightarrow$ Дискретизация времени (дает периодический непрерывный спектр, ДВПФ) $\rightarrow$ Усечение по времени до длины $N$ и Дискретизация спектра (дает дискретный спектр, ДПФ).

\section{Алгоритм быстрого преобразования Фурье (БПФ)}

Прямое вычисление ДПФ требует $O(N^2)$ комплексных умножений и сложений. Алгоритм Быстрого Преобразования Фурье (БПФ, FFT) позволяет сократить вычислительную сложность до $O(N \log_2 N)$. Рассмотрим исторически первый и наиболее популярный алгоритм \textbf{Кули-Тьюки по основанию 2 с прореживанием по времени} (Radix-2 DIT FFT).

\subsection{Идея алгоритма}
Пусть $N = 2^m$. Разделим исходную последовательность $x[n]$ на элементы с четными и нечетными индексами:
$$ X[k] = \sum_{n=0}^{N-1} x[n] W_N^{kn} = \sum_{r=0}^{N/2-1} x[2r] W_N^{2rk} + \sum_{r=0}^{N/2-1} x[2r+1] W_N^{(2r+1)k}. $$
Учитывая, что $W_N^{2} = e^{-j \frac{4\pi}{N}} = e^{-j \frac{2\pi}{N/2}} = W_{N/2}$, получаем:
$$ X[k] = \sum_{r=0}^{N/2-1} x_{even}[r] W_{N/2}^{rk} + W_N^k \sum_{r=0}^{N/2-1} x_{odd}[r] W_{N/2}^{rk} = E[k] + W_N^k O[k], $$
где $E[k]$ и $O[k]$ — ДПФ длины $N/2$ от четных и нечетных отсчетов соответственно.
Поскольку $E[k]$ и $O[k]$ периодичны с периодом $N/2$, а $W_N^{k+N/2} = -W_N^k$, вторая половина спектра вычисляется по формуле:
$$ X\left[k + \frac{N}{2}\right] = E[k] - W_N^k O[k]. $$
Эта пара уравнений называется \textbf{«бабочкой»} (Butterfly). Она позволяет рекурсивно сводить задачу длины $N$ к двум задачам длины $N/2$.

\subsection{Шаги алгоритма БПФ}
\begin{enumerate}
    \item \textbf{Бит-реверсивная перестановка}: Элементы массива $x[n]$ переставляются в порядке, соответствующем инверсии битов их индексов (например, индекс 1 для $N=8$ в двоичном виде `001` становится `100`, то есть 4). Это позволяет проводить вычисления "на месте" (in-place).
    \item \textbf{Базовый случай}: Рассматриваются пары элементов (ДПФ длины 2). Выполняется операция "бабочка": сложение и вычитание соседних элементов.
    \item \textbf{Слияние подзадач}: На каждом следующем этапе $s$ (от $1$ до $\log_2 N$) длина объединяемых блоков удваивается ($L = 2^s$). Выполняются операции "бабочка" с умножением на соответствующие поворачивающие множители $W_L^k$.
    \item \textbf{Окончание}: Процесс завершается, когда длина блока достигает $N$.
\end{enumerate}

\subsection{Пример работы БПФ для $N=4$}
Пусть задан сигнал $x = [1, 2, 3, 4]$. Вычислим его ДПФ через БПФ. $N=4, W_4 = e^{-j\pi/2} = -j$.

\textbf{Шаг 1. Перестановка.} Индексы: $0(00), 1(01), 2(10), 3(11)$.
Бит-реверс: $00(0), 10(2), 01(1), 11(3)$.
Переставленный массив: $x' = [x[0], x[2], x[1], x[3]] = [1, 3, 2, 4]$.

\textbf{Шаг 2. БПФ длины 2 ($s=1$, блоки по 2).} $W_2^0 = 1$.
Для блока $[1, 3]$ (четные исходного):
$E[0] = 1 + 3 = 4$, $E[1] = 1 - 3 = -2$.
Для блока $[2, 4]$ (нечетные исходного):
$O[0] = 2 + 4 = 6$, $O[1] = 2 - 4 = -2$.
Промежуточные результаты: $E = [4, -2]$, $O = [6, -2]$.

\textbf{Шаг 3. БПФ длины 4 ($s=2$, объединяем $E$ и $O$).} Поворачивающие множители: $W_4^0 = 1, W_4^1 = -j$.
Формулы: $X[k] = E[k] + W_4^k O[k]$, $X[k+2] = E[k] - W_4^k O[k]$ для $k=0, 1$.

Для $k=0$:
$X[0] = E[0] + W_4^0 O[0] = 4 + 1 \cdot 6 = 10$.
$X[2] = E[0] - W_4^0 O[0] = 4 - 1 \cdot 6 = -2$.

Для $k=1$:
$X[1] = E[1] + W_4^1 O[1] = -2 + (-j)(-2) = -2 + 2j$.
$X[3] = E[1] - W_4^1 O[1] = -2 - (-j)(-2) = -2 - 2j$.

\textbf{Ответ:} $X = [10, -2+2j, -2, -2-2j]$.
Алгоритм сократил количество комплексных умножений благодаря переиспользованию вычислений в "бабочках".

\end{document}
